\chapter{2003年京蒙皖卷（春文）}
\section{单选题}
\begin{problem}
设 \(a\)，\(b\)，\(c\)，\(d\in\R\) 且 \(a>b\)，\(c>d\)，且下列结论中正确的是（\quad）

\choices
    {\(a+c>b+d\)}
    {\(a-c>b-d\)}
    {\(ac>bd\)}
    {\(\dfrac{a}{d}>\dfrac{b}{c}\)}
\end{problem}

\begin{problem}
设 \(M\) 和 \(m\) 分别表示函数 \(y=\dfrac{1}{3}\cos x-1\) 的最大值和最小值，则 \(M+m\) 等于（\quad）

\choices
    {\(\dfrac{2}{3}\)}
    {\(-\dfrac{2}{3}\)}
    {\(-\dfrac{4}{3}\)}
    {\(-2\)}
\end{problem}

\begin{problem}
若 \(f(x)=\dfrac{x-1}{x}\)，则方程 \(f(4x)=x\) 的根是（\quad）

\choices
    {\(-2\)}
    {\(2\)}
    {\(-\dfrac{1}{2}\)}
    {\(\dfrac{1}{2}\)}
\end{problem}

\begin{problem}
若集合 \(M=\{y\mid y=2^{-x}\}\)，\(P=\{y\mid y=\sqrt{x-1}\}\)，则 \(M\cap P=\)（\quad）

\choices
    {\(\{y\mid y>1\}\)}
    {\(\{y\mid y\geqslant1\}\)}
    {\(\{y\mid y>0\}\)}
    {\(\{y\mid y\geqslant0\}\)}
\end{problem}

\begin{problem}
若 \(A\)，\(B\)，\(C\) 是 \(\triangle ABC\) 的三个内角，且 \(A<B<C\)（\(C\neq\dfrac{\pi}{2}\)），则下列结论中正确的是（\quad）

\choices
    {\(\tan A<\tan C\)}
    {\(\cot A<\cot C\)}
    {\(\sin A<\sin C\)}
    {\(\cos A<\cos C\)}
\end{problem}

\begin{problem}
在等差数列 \(\{a_n\}\) 中，已知 \(a_1+a_2+a_3+a_4+a_5=20\)，那么 \(a_3\) 等于（\quad）

\choices
    {\(4\)}
    {\(5\)}
    {\(6\)}
    {\(7\)}
\end{problem}

\begin{problem}
设复数 \(z_1=-1+\i\)，\(z_2=\dfrac{1}{2}+\dfrac{\sqrt{3}}{2}\i\)，则 \(\arg\dfrac{z_1}{z_2}=\)（\quad）

\choices
    {\(-\dfrac{5\pi}{12}\)}
    {\(\dfrac{5\pi}{12}\)}
    {\(\dfrac{7\pi}{12}\)}
    {\(\dfrac{13\pi}{12}\)}
\end{problem}

\begin{problem}
函数 \(f(x)=\lvert x\rvert\) 和 \(g(x)=x(2-x)\) 的递增区间依次是（\quad）

\choices
    {\((-\infty,0]\)，\((-\infty,1]\)}
    {\((-\infty,0]\)，\([1,+\infty)\)}
    {\([0,+\infty)\)，\((-\infty,1]\)}
    {\([0,+\infty)\)，\([1,+\infty)\)}
\end{problem}

\begin{problem}
在同一坐标系中，方程 \(a^2x^2+b^2y^2=1\) 与 \(ax+by^2=0\)（\(a>b>0\)）的曲线大致是（\quad）

\choices
    {\choicebitmap[width=0.15\paperwidth]{img/2003/jing_meng_wan_liberal_spring/q9_optA.png}}
    {\choicebitmap[width=0.15\paperwidth]{img/2003/jing_meng_wan_liberal_spring/q9_optB.png}}
    {\choicebitmap[width=0.15\paperwidth]{img/2003/jing_meng_wan_liberal_spring/q9_optC.png}}
    {\choicebitmap[width=0.15\paperwidth]{img/2003/jing_meng_wan_liberal_spring/q9_optD.png}}
\end{problem}

\begin{problem}
某班新年联欢会原定的 \(5\) 个节目已排成节目单，开演前又增加了两个新节目. 如果将这两个节目插入原节目单中，且两个新节目不相邻，那么不同插法的种数为（\quad）

\choices
    {\(6\)}
    {\(12\)}
    {\(15\)}
    {\(30\)}
\end{problem}

\begin{problem}
如图，在正三角形 \(ABC\) 中，\(D\)，\(E\)，\(F\) 分别为各边的中点，\(G\)，\(H\)，\(I\)，\(J\) 分别为 \(AF\)，\(AD\)，\(BE\)，\(DE\) 的中点. 将 \(\triangle ABC\) 沿 \(DE\)，\(EF\)，\(DF\) 折成三棱锥以后，\(GH\) 与 \(IJ\) 所成角的度数为（\quad）

\bitmapfigure[float=inline,max width=0.42\linewidth]{img/2003/jing_meng_wan_liberal_spring/q11_fig1.png}

\choices
    {\(90^{\circ}\)}
    {\(60^{\circ}\)}
    {\(45^{\circ}\)}
    {\(0^{\circ}\)}
\end{problem}

\begin{problem}
已知直线 \(ax+by+c=0\)（\(abc\neq0\)）与圆 \(x^2+y^2=1\) 相切，则三条边长分别为 \(\lvert a\rvert\)，\(\lvert b\rvert\)，\(\lvert c\rvert\) 的三角形（\quad）

\choices
    {是锐角三角形}
    {是直角三角形}
    {是钝角三角形}
    {不存在}
\end{problem}

\section{填空题}
\begin{problem}
函数 \(y=\sin 2x+1\) 的最小正周期为\fillinblank{}.
\end{problem}

\begin{problem}
如图，一个底面半径为 \(R\) 的圆柱形量杯中装有适量的水. 若放入一个半径为 \(r\) 的实心铁球，水面高度恰好升高 \(r\)，则 \(\dfrac{R}{r}=\)\fillinblank{}.

\bitmapfigure[max width=0.52\linewidth]{img/2003/jing_meng_wan_liberal_spring/q14_fig1.png}
\end{problem}

\begin{problem}
在某报《自测健康状况》的报道中，自测血压结果与相应年龄的统计数据如下表. 观察表中数据的特点，用适当的数填入表中空白内.

\begin{center}
\begin{tabular}{c|cccccccc}
年龄（岁） & \(30\) & \(35\) & \(40\) & \(45\) & \(50\) & \(55\) & \(60\) & \(65\) \\
\hline
收缩压（水银柱/毫米） & \(110\) & \(115\) & \(120\) & \(125\) & \(130\) & \(135\) &\fillinblank{} & \(145\) \\
舒张压（水银柱/毫米） & \(70\) & \(73\) & \(75\) & \(78\) & \(80\) & \(83\) &\fillinblank{} & \(88\)
\end{tabular}
\end{center}
\end{problem}

\begin{problem}
已知 \(F_1\)，\(F_2\) 分别为椭圆 \(\dfrac{x^2}{a^2}+\dfrac{y^2}{b^2}=1\) 的左，右焦点，点 \(P\) 在椭圆上，\(\triangle POF_2\) 是面积为 \(\sqrt{3}\) 的正三角形，则 \(b^2\) 的值是\fillinblank{}.
\end{problem}

\section{解答题}
\begin{problem}
解不等式：\(\log_{2}(x^2-x-2)>\log_{2}(2x-2)\).
\end{problem}

\begin{problem}
已知函数 \(f(x)=\dfrac{6\cos^4x-5\cos^2x+1}{\cos 2x}\).
求 \(f(x)\) 的定义域，判断它的奇偶性，并求其值域.
\end{problem}

\begin{problem}
如图，\(ABCD-A_1B_1C_1D_1\) 是正四棱柱，侧棱长为 \(1\)，底面边长为 \(2\)，\(E\) 是棱 \(BC\) 的中点.

\begin{enumerate}
\item 求三棱锥 \(D_1-DBC\) 的体积；
\item 证明 \(BD_1\parallel\) 平面 \(C_1DE\)；
\item 求面 \(C_1DE\) 与面 \(CDE\) 所成二面角的正切值.
\end{enumerate}

\bitmapfigure[max width=0.54\linewidth]{img/2003/jing_meng_wan_liberal_spring/q19_fig1.png}
\end{problem}

\begin{problem}
设 \(A(-c,0)\)，\(B(c,0)\)（\(c>0\)）为两定点，动点 \(P\) 到 \(A\) 点的距离与到 \(B\) 点的距离的比为定值 \(a\)（\(a>0\)），求 \(P\) 点的轨迹.
\end{problem}

\begin{problem}
某租赁公司拥有汽车 \(100\) 辆. 当每辆车的月租金为 \(3000\text{ 元}\) 时，可全部租出. 当每辆车的月租金每增加 \(50\text{ 元}\) 时，未租出的车将会增加一辆. 租出的车每辆每月需要维护费 \(150\text{ 元}\)，未租出的车每辆每月需要维护费 \(50\text{ 元}\).

\begin{enumerate}
\item 当每辆车的月租金定为 \(3600\text{ 元}\) 时，能租出多少辆车？
\item 当每辆车的月租金定为多少元时，租赁公司的月收益最大，最大月收益是多少？
\end{enumerate}
\end{problem}

\begin{problem}
如图，在边长为 \(l\) 的等边 \(\triangle ABC\) 中，圆 \(O_1\) 为 \(\triangle ABC\) 的内切圆，圆 \(O_2\) 与圆 \(O_1\) 外切，且与 \(AB\)，\(BC\) 相切，\(\cdots\)，圆 \(O_{n+1}\) 与圆 \(O_n\) 外切，且与 \(AB\)，\(BC\) 相切，如此无限继续下去. 记圆 \(O_n\) 的面积为 \(a_n\)（\(n\in\N\)）.

\begin{enumerate}
\item 证明 \(\{a_n\}\) 是等比数列；
\item 求 \(\displaystyle\lim_{n\to\infty}(a_1+a_2+\cdots+a_n)\) 的值.
\end{enumerate}

\bitmapfigure[max width=0.48\linewidth]{img/2003/jing_meng_wan_liberal_spring/q22_fig1.png}
\end{problem}
